\chapter*{Conclusion and Future Work}
\addcontentsline{toc}{chapter}{Conclusion and Future Work}

This study examined antimicrobial resistance prediction in \textit{Salmonella enterica} using curated AMR determinants derived from genomic data. The work combined cohort alignment, phenotype cleaning, genotype deduplication, class/subclass normalization, scope-controlled genotype-to-phenotype connection, and antibiotic-specific predictive modeling into one reproducible determinant-based workflow. The final report therefore contributes not only a set of prediction results, but also a clear data-processing methodology for turning curated surveillance-style resources into interpretable AMR model inputs.

The central scientific result is that the genotype-to-phenotype connection rule matters substantially. The project compared three connection scopes: \textit{strict}, which keeps only narrow lineage-matched determinants; \textit{broad}, which allows broader within-class determinant evidence and therefore captures cross-resistance; and \textit{all}, which uses the full determinant background observed in the isolate and therefore offers maximum coverage at the cost of greater co-occurrence leakage risk. The results show that \textit{all} is the strongest global default for predictive performance because it produces the highest average ROC-AUC and eliminates the zero-feature problem entirely. At the same time, \textit{broad} is the strongest biologically constrained alternative. It often performs nearly as well as \textit{all} and is especially useful when same-class resistance mechanisms need to be captured without opening the feature space to all detected determinants. \textit{Strict} remains valuable as a narrow biological reference and performs extremely well for some antibiotics, but it is too sparse to serve as a universal default.

The comparison with rule-based baselines clarifies why machine learning is useful in this setting. Direct determinant-based rules are already strong for some antibiotics whose phenotypes are tightly linked to specific genes or mutations. However, once broader within-class evidence must be considered, direct rules become too coarse. Machine learning preserves the benefit of wider scope while combining determinant evidence more selectively, which improves precision without sacrificing recall as severely as broad rule-based inference. Another major conclusion is that feature engineering decisions in this project are not merely technical. The choice of phenotype thresholds, the way duplicated genotype detections are collapsed, the handling of combined class/subclass annotations, the manual resolution of rare but meaningful edge cases, and the final scope rule all directly shape the biological meaning of the model input. The pipeline is therefore best understood as a bioinformatics workflow with a supervised learning stage, rather than as a classifier benchmark alone.

Feature-importance analysis further supports the main interpretation. Under narrow scopes, top-ranked features often align closely with expected resistance mechanisms, while the widest scope reveals a mixture of biologically plausible signals and co-occurring determinants from other resistance contexts. This reinforces the conclusion that predictive gain and interpretability do not increase together automatically; they must be balanced through explicit scope design. Several directions remain for future work. The most immediate extensions are external validation on an independent cohort, MIC prediction instead of binary phenotype prediction, and richer biological interpretation of top-ranked determinants at the antibiotic level. A broader extension would be to compare curated determinant-based workflows against wider resistome or pan-genome feature representations while keeping the same antibiotic-specific evaluation philosophy. These directions would test how much additional predictive value can be gained without losing the interpretability that motivated the determinant-based design in the first place.
